\documentclass[10pt,a4paper]{amsart}
\usepackage[margin=19mm]{geometry}
\usepackage{amsmath,amssymb,mathtools}
\usepackage[scr=rsfs,cal=euler]{mathalfa}
\usepackage{xfrac}
\usepackage[unicode,hidelinks,bookmarks=false]{hyperref}
\newcommand{\ie}{i.e.\ }
\newcommand{\cf}{cf.\ }
\newcommand{\resp}{respectively\ }
\newcommand{\Subsection}{\subsection}
\providecommand{\nobreakdash}{\nobreak\hbox{-}}
\setlength{\emergencystretch}{2em}
\multlinegap0pt
\usepackage{amssymb}
\usepackage{xcolor}
\usepackage{dsfont}
\usepackage{tikz}
\usepackage{mathtools}
\newtheorem{theorem}{Theorem}
\newtheorem{proposition}{Proposition}
\newtheorem{lemma}{Lemma}
\def\E{\mathbb{E}}
\newcommand{\R}{\mathbb{R}}
\def\T{\mathbb{T}}
\newcommand{\cjd}{\rangle}
\newcommand{\cjg}{\langle}
\def\eps{\varepsilon}
\renewcommand{\geq}{\;\geqslant\;}                   % redef. of > or =
\newcommand{\la}{\lambda}
\renewcommand{\leq}{\;\leqslant\;}                   % redef. of < or =
\newcommand{\mc}{\mathcal}
\title{2d Sinh-Gordon model on the infinite cylinder}
\author{Colin Guillarmou, Trishen S. Gunaratnam, Vincent Vargas}
\date{Source: arXiv:2405.04076v4 (CC BY 4.0)}
\begin{document}
\maketitle

\noindent\emph{Source and licence.} 2d Sinh-Gordon model on the infinite cylinder. Version arXiv:2405.04076v4 (CC BY 4.0), distributed by arXiv under the Creative Commons Attribution 4.0 International licence, http://creativecommons.org/licenses/by/4.0/, as recorded in the arXiv licence metadata for this exact version and shown on the abstract page https://arxiv.org/abs/2405.04076v4. Full licence notice: \texttt{licenses/arxiv-2405.04076-CC-BY-4.0.txt}.\par\smallskip

\noindent\textit{Editorial scope.} Counted unit: the article's theorem correlations (source0.tex lines 452--502). The article's complete original proof of it is delivered with it (source0.tex lines 1755--1818, one \texttt{proof} environment of 64 lines, which the article places away from the statement). No mathematics is written, repaired or re-derived: the statement and the proof are the article's own lines, in the article's own notation, and no retained line is substituted or re-flowed.  Retained uncounted as the source-local context the proof needs: lines 362--373; lines 1443--1456; lines 1639--1639; lines 1653--1657; lines 1673--1677; lines 1687--1704; lines 1695--1701.  Nothing was omitted from any retained span, so the package carries no omission record.  No retained line contains a citation, so the wrapper carries no bibliography and no bibliography entry.  No retained line contains a figure environment, an image-inclusion command or drawing code, so no image is delivered and the package declares no asset dependency.  Editorial formatting only, in addition: an AMS \texttt{amsart} preamble that restates the article's own declarations and macros used by the retained lines, namely the environments \texttt{theorem}, \texttt{proposition}, \texttt{lemma} on separate counters, together with the standard packages those lines need.  The retained environments print in this document as Theorem 1, Proposition 1, Lemma 1 in order of appearance, whereas the article prints the same items with the numbering of its own full document.  The complete native source of the article as distributed in the arXiv e-print for this exact version is bundled unchanged as \texttt{sources/158767/source0.tex}.
\begin{theorem}\label{thm:spectrum_H}
Let $\mu > 0$ and $\gamma\in(0,2)$. The Hamiltonian ${\bf H}$ generating the semigroup ${\bf T}_t$ is a self-adjoint, positive operator acting on a dense domain $\mathcal{D}(\mathbf{H}) \subset \mathcal{H}$. Furthermore, the following properties hold:
\begin{enumerate}
\item The operator $\mathbf{H}$ has discrete spectrum $(\lambda_j)_{j \geq 0}\subset \R_+$  with complete basis of normalised eigenfunctions $(\psi_j)_{j \geq 0}$,  the smallest eigenvalue $\lambda_0$ is simple and strictly positive, and there is  a unique associated normalised eigenfunction $\psi_0\in \mc{H}$, which is positive $dc\otimes d\mathbb P_\T$-almost everywhere. We call $\psi_0$ the ground state.
\item For every $j \geq 0$ and $N \geq 0$, the eigenfunctions $\psi_j$ belong to the weighted $L^p$ spaces $e^{-N|c|}L^p(H^{-s}(\T),\mu_0)$.
\item  One has the following diagonalization result
\begin{equation}\label{introdiagonalize}
e^{-t{\bf H}}f(c+\varphi)= \sum_{j=0}^{\infty} e^{-\lambda_j t}\psi_j(c+\varphi)
\langle f, \psi_j \rangle_{\mc{H}}, \qquad \forall f \in \mathcal{H}.
\end{equation}
\end{enumerate}
\end{theorem}

\begin{proposition} \label{prop: sinh expectation}
Let $F \in \mathcal{F}$ with support $[t_1,t_2]$ where $-\infty < t_1 \leq t_2 < \infty$. Then 
\[\lim_{T\to \infty} \langle F \rangle_{\mc{C}_{1,T}}
=\cjg F\cjd_{\mc{C}}.\]
Moreover,  for each sequence $T_n\to \infty$, there is a subsequence $T_{n_k}$ such for $\mu_0$-almost every $c+\varphi\in H^{-s}(\T)$ (with $s>0$),
\begin{equation}\label{ae_convergence}
\lim_{k\to \infty} \frac{1}{(e^{-T_{n_k}{\bf H}}1)(c+\varphi)}  \E_\varphi[F (c+B_{\bullet+T_{n_k}}+\varphi_{\bullet+T_{n_k}}) \,   e^{- \mu\sum_{\sigma = \pm 1} 
e^{\sigma\gamma c} M_\gamma^\sigma(\mc{C}_{[0,2T_{n_k}]})}]
=\cjg F\cjd_{\mc{C}}.	
\end{equation}
Finally, the first eigenfunction and eigenvalues can be obtained by the expressions 
\[ \la_0=-\frac{1}{2}\lim_{T\to +\infty}\frac{1}{T}\log Z_{\mc{C}_{1,T}}, \quad \Big(\int \psi_0d\mu_0\Big) \psi_0= \lim_{T\to \infty}e^{\la_0T}e^{-T{\bf H}}1,\]
where the second limit is with the topology of $\mc{H}$.
\end{proposition}

\subsection{Vertex correlations} \label{subsec:_vertex_correlations}

\begin{equation}\label{defmcI}
\mathcal I
=
\{ (\alpha_{ij}, (t_i, \theta_{ij})) \,|\,  i,j\in [1,N],  0 < t_1 < \dots < t_N<t \}.	
\end{equation}

\begin{equation}\label{def_of_h}
h_\eps(\theta)
:=
\sum_{i,j} \alpha_{ij}\mathbb E[ \varphi(\theta) \varphi^\eps_{t_i}(\theta_{ij})], \quad h(\theta): = \sum_{i,j} \alpha_{ij}\mathbb E[ \varphi(\theta) \varphi_{t_i}(\theta_{ij})].
\end{equation}

\begin{lemma}\label{lem: vertex correlations}
For every $\beta > 0$, there exists $C>0$ such that
\begin{equation} \label{eq: vc unif l2}
\limsup_{\eps \rightarrow 0}\int \mathbb{E}[G_{\beta,\eps}(c+\varphi)^2] e^{-\beta |c|}dc
\leq
C. 
\end{equation}
Consequently, with the notation \eqref{defmcI}, one has 
\begin{align}\label{eq: vc limit}
\lim_{\eps \rightarrow 0}\Big\cjg \prod_{(\alpha,z)\in \mc{I}}V^\eps_{\alpha}(z)\Big\cjd_{\mc{C}}	
&=
e^{\lambda_0 t+\frac{1}{2}\sum_{i,j}\alpha_{ij}^2t_i}  \int_{\R} e^{\sum_{i,j}\alpha_{ij} c}\mathbb{E}\big[\psi_0\big(c+B_t+\sum_{i,j} \alpha_{ij}t_i+ \varphi_t + Ph (t,\cdot))\psi_0(c+\varphi+h) 
\\
& \qquad \qquad \times e^{-\mu \sum_{\sigma = \pm 1} e^{\sigma\gamma c} \int_{\mathcal{C}_{[0,t]}} \prod_{i,j} |e^{-s+i\theta}-e^{-t_i+i\theta_{ij}}|^{-\gamma \sigma \alpha_{ij}} M^\sigma_\gamma(dsd\theta)}\big ]dc
\end{align}
and the limiting correlation is nontrivial if and only if $\mathcal I$ is $\gamma$-admissible.

\end{lemma}

\begin{align}
\lim_{\eps \rightarrow 0}\Big\cjg \prod_{(\alpha,z)\in \mc{I}}V^\eps_{\alpha}(z)\Big\cjd_{\mc{C}}	
&=
e^{\lambda_0 t+\frac{1}{2}\sum_{i,j}\alpha_{ij}^2t_i}  \int_{\R} e^{\sum_{i,j}\alpha_{ij} c}\mathbb{E}\big[\psi_0\big(c+B_t+\sum_{i,j} \alpha_{ij}t_i+ \varphi_t + Ph (t,\cdot))\psi_0(c+\varphi+h) 
\\
& \qquad \qquad \times e^{-\mu \sum_{\sigma = \pm 1} e^{\sigma\gamma c} \int_{\mathcal{C}_{[0,t]}} \prod_{i,j} |e^{-s+i\theta}-e^{-t_i+i\theta_{ij}}|^{-\gamma \sigma \alpha_{ij}} M^\sigma_\gamma(dsd\theta)}\big ]dc
\end{align}

\begin{theorem} \label{thm: sinh correlations}
Let $\mu>0$, $\gamma\in (0,2)$, $R>0$ and fix a set $\mathcal{I}=\bigsqcup_{i=1}^n\mc{I}_{t_i}$ of insertions with $\mathcal{I}_{t_i}= \cup_{j=1}^{l_i} \lbrace   (\alpha_{ij},t_i, \theta_{ij}) \rbrace$. The vertex correlations are defined via the formula
\[\Big\cjg \prod_{(\alpha,z)\in \mc{I}}V_{\alpha}(z)\Big\cjd_{\mc{C}_R,\mu} :=\lim_{\eps\to 0}
\Big\langle \prod_{(\alpha, (t, \theta)) \in \mathcal{I}}\eps^{\alpha^2/2}e^{\alpha (c+ B_{t/R} +  \varphi^{R,\eps}(t,\theta)) } \Big\rangle_{\mc{C}_R,\mu}\, ,
\]
where the limit exists for any insertion set $\mathcal I$ and is nontrivial if and only if $\mathcal{I}$ is $\gamma$-admissible. In addition, $\gamma$-admissible vertex correlations obey the following relations:
\vspace{3mm}

\noindent 1) Let $\mathcal{I}$ be $\gamma$-admissible insertions on $\mc{C}_R$ and let $\mathcal{I}_{1/R} = \{ (\alpha, z/R) \,|\, (\alpha,z) \in \mathcal{I} \}$. Then we have the following scaling relation: 
\begin{equation}\label{scale invariance}
\Big\cjg \prod_{(\alpha,z)\in \mc{I}}V_{\alpha}(z)\Big\cjd_{\mc{C}_R,\mu} 
=
R^{\sum_{(\alpha,z) \in \mathcal{I}} \frac{\alpha^2}{2} }
\Big\cjg \prod_{(\alpha,z)\in \mc{I}_{1/R}}V_{\alpha}(z)\Big\cjd_{\mc{C},\mu_R}.
\end{equation}
with $\mu_R=R^{\gamma Q}\mu$.
 In particular, the one-point function $\langle V_\alpha(0) \rangle_{\mathcal C_R}:=\langle e^{\alpha \varphi(0)} \rangle_{\mc{C}_R,\mu}$ exists, is nontrivial for $|\alpha|<Q$, and obeys the scaling relation
\begin{equation} \label{one-pt}
\langle e^{\alpha \varphi(0)} \rangle_{\mc{C}_R,\mu}
=
R^{\alpha^2/2} \langle e^{\alpha \varphi(0)} \rangle_{\mc{C},\mu_R}.	
\end{equation}
\vspace{2mm}

\noindent 2)
One has the following formula for the correlations when $R=1$; if $t_i >0$ for all $i$, for any $t>\max(t_1,\dots,t_n)$ 
\begin{align}
\begin{split}
 \Big\cjg \prod_{(\alpha,z)\in \mc{I}}V_{\alpha}(z)\Big\cjd_{\mc{C},\mu}   & =  e^{\lambda_0 t + \frac{1}{2}\sum_{i,j}\alpha_{ij}^2t_i} \int_\R e^{\sum_{i,j}\alpha_{ij} c}\mathbb{E}\big[\psi_0\big(c+B_t+\sum_{i,j} \alpha_{ij}t_i+ \varphi_t + Ph (t,\cdot))\psi_0(c+\varphi+h) 
\\
& \qquad \qquad \times e^{-\mu \sum_{\sigma = \pm 1} e^{\sigma\gamma c} \int_{[0,t]\times \T} \prod_{i,j} |e^{-s+i\theta}-e^{-t_i+i\theta_{ij}}|^{-\gamma \sigma \alpha_{ij}} M^\sigma_\gamma(dsd\theta)}\big ]dc
\end{split}
\end{align}
where $Ph(t,\theta) = \sum_{i,j} \alpha_{ij}\mathbb E[ \varphi_t(\theta) \varphi_{t_i}(\theta_{ij})]$ is a harmonic function with boundary value $h(\theta)= \sum_{i,j} \alpha_{ij}\mathbb E[ \varphi(\theta) \varphi_{t_i}(\theta_{ij})]$ at $t=0$ and decaying to $0$ as $t\to \infty$.
If $\mathcal{I}=\{(\alpha,0)\}$ with $\alpha\in (-Q,Q)$, then  \begin{equation} 
\label{vertex explicit}
\cjg V_{\alpha}(0)\cjd_{\mc{C},\mu} = \|e^{\alpha c/2} \psi_{0}(\cdot+h)\|^2_\mc{H}
\end{equation}
with $h(\theta)=-\alpha \log|e^{i\theta}-1|$.

\vspace{3mm}

\noindent 3) For every $t \in \mathbb R$, $\theta \in \mathbb T_R$, and $\alpha \in (-Q,Q)$, we have that the corresponding two-point vertex correlations\footnote{Similar statements hold for $F,G:C(\mathbb R, H^{-s}(\mathbb T_R))\to \R$ that are bounded, measurable, and of compact support. See the proof of Theorem \ref{thm: sinh correlations} in Section \ref{subsec:_vertex_correlations}.} are positively correlated:
\begin{equation}
{\rm Cov}_{\mathcal C_R}(V_\alpha(0,\theta), V_\alpha(t,\theta)):= \langle V_\alpha(0,\theta) V_\alpha(t,\theta) \rangle_{\mathcal C_R} - \langle V_\alpha(0,\theta) \rangle_{\mathcal C_R}\langle  V_\alpha(t,\theta) \rangle_{\mathcal C_R}\geq 0. 	
\end{equation}
Furthermore, we have exponential decay of the two-point vertex correlation functions: for every $\theta_1,\theta_2 \in \mathbb T_R$ and every $\alpha_1,\alpha_2 \in (-Q,Q)$, there exists $C_{R}(\alpha_1,\theta_1; \alpha_2,\theta_2)>0$ such that, for every $|t| > 1$,
\begin{equation}
|{\rm Cov}_{\mathcal C_R}(V_{\alpha_1}(0,\theta_1), V_{\alpha_2}(t,\theta_2))| \leq e^{-(\lambda_1^{\mu_R}-\lambda_0^{\mu_R})\frac{|t|}R} C_R(\alpha_1,\theta_1; \alpha_2,\theta_2).	
\end{equation}
\end{theorem}

\begin{proof}[Proof of Theorem \ref{thm: sinh correlations}]
For $R=1$, the convergence, nontriviality, and explicit formula of vertex correlations with arbitrary insertion set follows directly from Lemma \ref{lem: vertex correlations}. The generic case $R>0$, as we have argued previously, follows from the scaling relation for GMC. We observe that for the $1$-point function, if the insertion is at $(t_1,\theta_1)=(0,0)$ with weight $|\alpha|<Q$, we can let $t\to 0$ in the expression \eqref{eq: vc limit} and obtain (with $h(\theta)=-\alpha \log|e^{i\theta}-1|$ as in \eqref{def_of_h}):
\begin{equation}\label{1-point}
\cjg V_{\alpha}(0)\cjd_{\mc{C}}=  \|e^{\alpha c/2} \psi_0(\cdot + h)\|^2_{\mc{H}}.
\end{equation}


We now turn to the exponential decay of correlations for the truncated two-point function, or covariance, of the vertex correlations. We will fix $R=1$ for ease of notation; the case $R>0$ follows with straightforward modifications. As a preliminary step, we analyze covariances of bounded functions of compact support. Let $\eps > 0$. Using the same approach as in the proof of Proposition \ref{prop: sinh expectation}, for every $F_\eps \in \mathcal A$ bounded and supported on $[-\eps,\eps]$, define the map
\begin{equation}
\mc{B}_{F_\eps}:\mathcal H \rightarrow \mathcal H, \quad (\mc{B}_{F_\eps} \psi)(c,\varphi) = \mathbb E_\varphi[F_\eps(c+B_{\bullet +\eps} +\varphi_{\bullet +\eps}) \psi(c+B_{2\eps}+\varphi_{2\eps})e^{-\mu  \sum_{\sigma = \pm 1} e^{\sigma\gamma c}M_\gamma^\sigma([0,2\eps]\times\T)}].
\end{equation}
In particular, when $\eps = 0$ we have that $\mc{B}_{F_0}\psi = F_0 \psi  \in \mathcal {H}$. Given $F_\eps,G_\eps \in \mathcal A$ that are bounded and supported on $[-\eps,\eps]$, by the domain Markov property and tower property of conditional expectations, we have
\begin{equation}
\begin{split}
\langle F_\eps \cdot \tau_{t} G_\eps \rangle_{\mathcal C_1} =
%&  e^{\lambda_0 (t+2\eps)} \int_{\mathbb R} \mathbb E[\psi_0(c,\varphi)F_\eps(c+B_{\bullet +\eps} +\varphi_{\bullet +\eps})e^{-\mu  \sum_{\sigma = \pm} e^{\sigma\gamma c}M_\gamma^\sigma([0,2\eps]\times\T)} \, \cdot e^{-(t-2\eps)\mathbf H} \mc{B}_{G_\eps}(\psi_0)(c,\varphi)]\\
 e^{\lambda_0 (t+2\eps)} \cjg \mc{B}_{F_\eps}e^{-(t-2\eps){\bf H}}\mc{B}_{G_\eps}\psi_0,\psi_0\cjd_{\mc{H}},
\end{split}
\end{equation}
where we recall that $\tau_t G$ translates the support of $G$ by $+t$. 

It is straightforward to check that the operator $\mc{B}_{F_\eps}$ is self-adjoint (using that $F_\eps$ is real valued and the 
dynamics $B_s+\varphi_s$ is reversible with respect to $\mu_0$), thus the eigenfunction expansion in Theorem \ref{thm:spectrum_H} yields (for $t>2\eps)$: 
\[\begin{split}
\langle F_\eps \cdot \tau_{t} G_\eps \rangle_{\mathcal C_1} &= e^{\lambda_0(t+2\eps)} \sum_{j \geq 0} e^{-\lambda_j (t-2\eps)} \langle  \mc{B}_{F_\eps}\psi_0,\psi_j \rangle_{\mathcal H} \langle \mc{B}_{G_\eps}\psi_0,\psi_j\rangle_{\mathcal H}.
\end{split}\]
Hence since $\langle F_\eps \rangle_{\mathcal C_1}=e^{2\la_0\eps}\cjg \mc{B}_{F_\eps}\psi_0,\psi_0\cjd_{\mc{H}}$ and $\langle G_\eps \rangle_{\mathcal C_1}=e^{2\la_0\eps}\cjg \mc{B}_{G_\eps}\psi_0,\psi_0\cjd_{\mc{H}}$,
we have the representation
\begin{equation} \label{eq: covariances}
{\rm Cov}(F,\tau_t G) := \langle F_\eps \cdot \tau_tG_\eps \rangle_{\mathcal C_1} - \langle F_\eps \rangle_{\mathcal C_1} \langle \tau_t G_\eps \rangle_{\mathcal C_1} = e^{-(\lambda_1-\lambda_0)t} \mathcal R_t(F_\eps, G_\eps),
\end{equation}
where
\begin{equation}
\mathcal R_t(F_\eps, G_\eps):= e^{2(\lambda_1+\lambda_0)\eps} \sum_{j \geq 1} e^{-(\lambda_j - \lambda_1)(t-2\eps)} \langle  \mc{B}_{F_\eps}\psi_0,\psi_j \rangle_{\mathcal H} \langle \mc{B}_{G_\eps}\psi_0,\psi_j\rangle_{\mathcal H} 
\end{equation}
In particular, when $F_\eps = G_\eps$, the covariance is non-negative. 


Let us now turn to vertex correlations. Fix $t>0$ without loss of generality. Let $|\alpha_1|,|\alpha_2|<Q$ and $\theta_1,\theta_2 \in \mathbb T$. For each $K \in \mathbb N^*$, write
\begin{equation}
V_{\alpha}^\eps(s,\theta)_K:= \min(V_{\alpha}^\eps(s,\theta),K) \in \mathcal A,
\end{equation}
which is bounded and supported on $[-\eps, \eps]$. Hence, by \eqref{eq: covariances} we obtain that for every $\eps > 0$ such that $2\eps < t$ and for every $\delta > 0$,
\begin{equation}
{\rm Cov}\Big(V_{\alpha_1}^\eps(0,\theta_1)_K,  V_{\alpha_2}^\eps(t,\theta_2)_K  \Big) = e^{-(\lambda_1-\lambda_0)t} e^{(\lambda_1-\lambda_0)\delta} e^{-(\lambda_1-\lambda_0)\delta }\mathcal R_t(V_{\alpha_1}^\eps(0,\theta_1)_K,  V_{\alpha_2}^\eps(0,\theta_2)_K).
\end{equation}
We now choose $\eps$ and $\delta$ such that $2\eps < \delta < t$. Then by the Cauchy-Schwarz inequality, 
\begin{equation}
|e^{-(\lambda_1-\lambda_0)\delta }\mathcal R_t(V_{\alpha_1}^\eps(0,\theta_1)_K,  V_{\alpha_2}^\eps(0,\theta_2)_K)| \leq e^{-(\lambda_1-\lambda_0)\delta }\prod_{i=1,2} \mathcal R_\delta( 	V_{\alpha_i}^\eps(0,\theta_i)_K, V_{\alpha_i}^\eps(0,\theta_i)_K)^{1/2}.
\end{equation}
However, by \eqref{eq: covariances} and the fact that the covariances are positive when $F_\eps=G_\eps$, we obtain 
\begin{align*}
& \limsup_{\eps \rightarrow 0}\limsup_{K \rightarrow \infty}\Big|\prod_{i=1,2} \mathcal R_\delta( 	V_{\alpha_i}^\eps(0,\theta_i)_K, V_{\alpha_i}^\eps(0,\theta_i)_K)^{1/2}\Big|e^{-(\lambda_1-\lambda_0)\delta }	\\
& \leq \lim_{\eps \rightarrow 0} \lim_{K \rightarrow \infty} \prod_{i=1,2} {\rm Cov}\Big(V_{\alpha_i}^\eps(0,\theta_i)_K, V_{\alpha_i}^\eps(\delta,\theta_i)_K  \Big)^{1/2}.
\end{align*}
But, since the covariances of vertex correlations converge, we therefore have that
\begin{equation}
\Big|{\rm Cov}\Big(V_{\alpha_1}(0,\theta_1), V_{\alpha_2}(t,\theta_2)\Big)\Big| \leq e^{-(\lambda_1-\lambda_0)t} C(\alpha_1,\theta_1; \alpha_2,\theta_2),	
\end{equation}
where
\begin{equation}
C(\alpha_1,\theta_1;\alpha_2,\theta_2):= e^{(\lambda_1-\lambda_0)\delta} \prod_{i=1,2}{\rm Cov}\Big(V_{\alpha_i}(0,\theta_i), V_{\alpha_i}(\delta,\theta_i) \Big)^{1/2}. 	\qedhere
\end{equation}
\end{proof}

\end{document}
